% =====================================================================
%  ABSTRACT
% =====================================================================

\begin{resumo}[Abstract]
\begin{otherlanguage*}{english}

Virtual cycling has become a sport modality recognized by the Union
Cycliste Internationale (UCI), yet access to simulation platforms is
restricted by monthly subscription fees in foreign currency, high-performance
hardware requirements, and continuous internet dependence. This work presents
the software registration of Argus Cyclist, an open-source cross-platform
virtual cycling simulator licensed under GPLv3, developed in Go (Golang) with
a JavaScript frontend using the Wails framework for desktop and Capacitor.js
for Android mobile devices. The system employs a Clean Architecture with
Facade and Repository patterns, abstracting Bluetooth Low Energy (BLE)
connections for heart rate, power, and smart trainer control via GATT profiles
(\texttt{0x180D}, \texttt{0x1818}, \texttt{0x1826}). The physics engine
implements Martin's equation for speed calculation from GPX route gradients,
while the physiological analysis module computes advanced metrics such as
Normalized Power (NP), Training Stress Score (TSS), Performance Management
Chart (PMC), and Aerobic Decoupling. It distinguishes itself by incorporating
a local generative AI coaching assistant (Ollama) for creating personalized
workout plans in ZWO format, along with a bicycle component wear tracking
system. The software was registered at the Brazilian National Institute of
Industrial Property (INPI, Process BR512026003786-7) in its initial version,
having evolved significantly after the registration with new modules, training
modes, third-party integrations, and architecture improvements. Experimental
validation uses real cycling equipment to verify telemetry latency below
\SI{150}{\milli\second} and \texttt{.fit} file compliance with the Garmin
standard. The project aims to democratize scientific high-performance training
by eliminating economic and connectivity barriers.

\textbf{Keywords}: Virtual cycling. Open source. Physics simulation. Bluetooth
Low Energy. Physiological analysis. Artificial intelligence.

\end{otherlanguage*}
\end{resumo}
